\documentclass[12pt]{article}
% ----------------------------------------------------------------------
% Define external packages, language, margins, fonts, new commands 
% and colors
% ----------------------------------------------------------------------
\usepackage[utf8]{inputenc} % Codification
\usepackage[english]{babel} % Writing idiom

\usepackage[export]{adjustbox} % Align images
\usepackage{amsmath} %a Extra commands for math mode
\usepackage{amssymb} % Mathematical symbols
\usepackage{anysize} % Personalize margins
    \marginsize{2cm}{2cm}{2cm}{2cm} % {left}{right}{above}{below}
\usepackage{appendix} % Appendices
\usepackage{cancel} % Expression cancellation
\usepackage{caption} % Captions
    \captionsetup{labelfont={bf}}
\usepackage{cite} % Citations, like [1 - 3]
\usepackage{color} % Text coloring
\usepackage{fancyhdr} % Head note and footnote
    \pagestyle{fancy}
    \fancyhf{}
    \fancyhead[L]{
        \includegraphics[scale = 0.3]{Images/NovaFctHor.png}
    } % Left of Head note
    \fancyhead[R]{\footnotesize Tese} % Right of Head note
    \fancyfoot[L]{\footnotesize MEEC} % Left of Footnote
    \fancyfoot[R]{\thepage  } % Center of Footnote
    %\fancyfoot[R]{\footnotesize Degree} % Right of Footnote
    \renewcommand{\footrulewidth}{0.4pt} % Footnote rule
\usepackage{float} % Utilization of [H] in figures
\usepackage{graphicx} % Figures in LaTeX
\usepackage[colorlinks = true, plainpages = true, linkcolor = novablue, urlcolor = novablue, citecolor = novablue, anchorcolor = novablue]{hyperref}
\usepackage{indentfirst} % First paragraph
\usepackage[super]{nth} % Superscripts
\usepackage{siunitx} % SI units
\usepackage{subcaption} % Subfigures
\usepackage{titlesec} % Font
    \titleformat{\section}{\Large\bfseries}{\thesection}{1em}{}
    \titleformat{\subsection}{\large\bfseries}{\thesubsA_competitive_differential_game_between_an_unmanned_aerial_and_a_ground_vehicle_using_model_predictive_controlection}{1em}{}
    \titleformat{\subsubsection}{\normalsize\bfseries}{\thesubsubsection}{1em}{}
    %\fancyfoot[C]{\thepage}

%code
\usepackage[breakable]{tcolorbox}

% Code highlighting
\usepackage{listings}
\lstset{ 
    backgroundcolor=\color{cellbackground}, % background color for the code block
    basicstyle=\ttfamily,                   % font style
    breakatwhitespace=false,                % automatic breaks only at whitespace
    breaklines=true,                        % automatic line breaking
    captionpos=b,                           % caption position
    commentstyle=\color{gray},              % comment style
    escapeinside={\%*}{*)},                 % if you want to add LaTeX within your code
    keywordstyle=\color{blue},              % keyword style
    stringstyle=\color{dkgreen},            % string style
    frame=single,                           % adds a frame around the code
    rulecolor=\color{cellborder},           % border color
}

%---------TIKZ CONFIG---------%
%Block diagram
\usepackage{tikz}
\usetikzlibrary{shapes.geometric, arrows}
% Define custom color lightblue
\definecolor{lightblue}{RGB}{173, 216, 230}
\tikzstyle{block} = [draw, rectangle, 
    minimum height=2.5em, minimum width=3.5em]
\tikzstyle{arrow} = [thick,->,>=stealth]

\usepackage{xcolor}

% Define a custom light blue color
\definecolor{lightblue}{RGB}{173, 216, 230}

% Define styles with blue shades
\tikzstyle{startstop} = [rectangle, rounded corners, minimum width=3cm, minimum height=1cm, text centered, draw=black, fill=blue!30]
\tikzstyle{process}   = [rectangle, minimum width=3cm, minimum height=1cm, text centered, draw=black, fill=cyan!30]
\tikzstyle{decision}  = [diamond, minimum width=3cm, minimum height=1cm, text centered, draw=black, fill=lightblue]

%---------TIKZ CONFIG---------%

% Exact colors from NB
\definecolor{incolor}{HTML}{303F9F}
\definecolor{outcolor}{HTML}{D84315}
\definecolor{cellborder}{HTML}{CFCFCF}
\definecolor{cellbackground}{HTML}{F7F7F7}

%folders .tex
\usepackage{subfiles}

% Random text (not needed)
\usepackage{lipsum}
\usepackage{duckuments}

% Tables
\usepackage{tabularx} 
\usepackage{booktabs}
\usepackage{makecell}

% New and re-newcommands
\newcommand{\sen}{\operatorname{\sen}} % Sine function definition
\newcommand{\HRule}{\rule{\linewidth}{0.5mm}} % Specific rule definition
\renewcommand{\appendixpagename}{\LARGE Appendices}

% Colors

\definecolor{novablue}{RGB}{0, 101, 189}
\definecolor{dkgreen}{RGB}{52,129,65}

%%%%%%%%%%%%%%%%%%%%%%%%%%%%%%%%%%%%%%%%%%%%%%%%%%%%%%%%%%%%%%%%%%%%%%%%
%                                 Document                             %
%%%%%%%%%%%%%%%%%%%%%%%%%%%%%%%%%%%%%%%%%%%%%%%%%%%%%%%%%%%%%%%%%%%%%%%%
\begin{document}
% Add bibliography data and style
\bibliographystyle{IEEEtran}
% ----------------------------------------------------------------------
% Cover
% ----------------------------------------------------------------------
\begin{center}
    \begin{figure}
        \vspace{-1.0cm}
        \includegraphics[scale = 0.055, left]{Images/NovaFctVer.png} % Nova logo
        
    \end{figure}

    \mbox{}\\[2.0cm]
    \textsc{\Huge Mestrado em Engenharia Eletrotécnica e de Computadores}\\[2.5cm]
    \textsc{\LARGE Tese}\\[2.0cm]
    \HRule\\[0.4cm]
    {\large \bf {Estado de Arte}}\\[0.2cm]
    \HRule\\[1.5cm]
\end{center}

\begin{flushleft}
    \textbf{Author:}
\end{flushleft}

\begin{center}
    \begin{minipage}{0.5\textwidth}
        \begin{flushleft}
            Filipe Cavalheiro (70119)
        \end{flushleft}
    \end{minipage}%
    \begin{minipage}{0.5\textwidth}
        \begin{flushright}
            \href{mailto:fs.cavalheiro@campus.fct.unl.pt}{\texttt{fs.cavalheiro@campus.fct.unl.pt}}
        \end{flushright}
    \end{minipage}
\end{center}

\vspace{5cm}

\begin{center}
    \large \bf 2025/2026 -- \nth{1} Semester, DEEC
\end{center}

\thispagestyle{empty}

\setcounter{page}{0}

\newpage

\newpage

\tableofcontents % Generates the table of contents

\newpage

\listoffigures

\newpage

\textbf{Definition 1 (Enclosure Condition).} \textit{The evader is said to be enclosed in space by the pursuers if it lies inside the convex hull $\Omega$ formed by the pursuers, i.e. $p_e \in \Omega$.}\\

\textbf{Definition 2 (Capture Condition).} \textit{The evader is said to be captured if there exists at least one pursuer, whose distance from the evader is smaller than the capture radius denoted by $r_c$, i.e. $\exists j \in [m]$, such that $\| p_j - p_e \|_2 \leq r_c$. and the same is located above the evader}\\

The above definitions are inspired by: \cite{3D-Cooperative-Pursuit-with-Guaranteed-Enclosure-via-Robust-Model-Predictive-Control}

[TODO] insert the image of the drone used in project capture with its grabber to justify the need of the pursuer that captures to be above the evader. 

[IDEAS]
 The spherical shell of drones around the evader has a defined size that respects the safe distance for no accidental collision but is small enough to respect Definition 2. The position of pursuers in this shell is defined in a way that maximizes the distance between the pursuers. For example: for 2 pursuers the drones are located in opposite sides of the sphere being their max distance the diameter of the sphere. With 3 it would be a equilateral triangle in a plane. Etc, etc. Every time a new pursuer is added to the network the drones already around the evader move accordingly to this Definition [TODO] define it formally as Definition 3.

\newpage
\bibliography{references}

\end{document}